\begin{center}\small\begin{tabular}{lrrrrr}
\toprule
Policy & Revenue & Cost & Profit & Collected & Available \\ \midrule
Beta only & 52,625 & 5,844 & 46,781 & 37,262 & 38,787 \\
Private selective & 50,360 & 5,972 & 44,388 & 34,997 & 36,507 \\
Shared selective & 50,360 & 5,972 & 44,388 & 34,997 & 36,507 \\
Controlled Gym pilot & 48,092 & 6,221 & 41,871 & 32,729 & 34,104 \\
\bottomrule
\end{tabular}\end{center}
All entries are mean simulated USD over 30 days. Beta's operating profit is 46,781.09, with conditional seed interval [45,409.91, 48,152.27]. Its A/B/C costs are 3,719.76, 2,051.60 and 72.80. Mean collections are 37,262.50; receivables are 15,362.75. Cash of 43,879.23 includes the separate 12,000 seed capital. Reserves are 4,631.26; unpaid obligations are 460.88. Available funds are 38,787.08; available earned surplus is 26,787.08. These distinct quantities reconcile in every retained run.

The primary optional policies reject v1 and preserve cash instead of choosing Gym training. They still pay for assessing the candidate and lose delivery capacity while doing so. Their mean profit is USD 2,393.30 below Beta. A controlled Gym-pilot arm deliberately pays for development so that the experiment can test v2; it is not the endogenous optimizer's selected action. That arm rejects v2 at the unrestricted envelope and loses USD 4,910.24 relative to Beta. Rejecting the investment does not erase the cost of discovering that it was unattractive.\par In the predeclared 40-call fitting-envelope scenario, the controlled Gym pilot has mean conditional profit USD 10,599.20, versus Beta's 6,160.18. Its paired increment is 4,439.01, interval [3,957.35, 4,920.68]. Response-model contracts still pay 384 additional acceptance queries in both arms. This narrow-envelope economic result is separate from the robust-plan supplement and from the primary 96-call scenario. The pilot uses a point-forecast investment rule; the existing lifecycle instead rejects the candidate under its conservative uncertainty rule. The pilot is a controlled counterfactual, not an authorized shared-production outcome passing that gate. It depends on assumed buyer behavior, scarcity, prices and reuse, and does not establish net value from sharing rather than private use.
